\documentclass[10pt,a4paper]{amsart}
\usepackage[margin=19mm]{geometry}
\usepackage{amsmath,amssymb,mathtools}
\usepackage[scr=rsfs,cal=euler]{mathalfa}
\usepackage{xfrac}
\usepackage[unicode,hidelinks,bookmarks=false]{hyperref}
\newcommand{\ie}{i.e.\ }
\newcommand{\cf}{cf.\ }
\newcommand{\resp}{respectively\ }
\newcommand{\Subsection}{\subsection}
\providecommand{\nobreakdash}{\nobreak\hbox{-}}
\setlength{\emergencystretch}{2em}
\multlinegap0pt
\usepackage{amssymb}
\usepackage{bm}
\newtheorem{lemma}{Lemma}
\newtheorem{prop}{Proposition}
\title{Envelopes and evolutes}
\author{Ragni Piene}
\date{Source: arXiv:2508.08145v2 (CC BY 4.0)}
\begin{document}
\maketitle

\noindent\emph{Source and licence.} Envelopes and evolutes. Version arXiv:2508.08145v2 (CC BY 4.0), distributed by arXiv under the Creative Commons Attribution 4.0 International licence, http://creativecommons.org/licenses/by/4.0/, as recorded in the arXiv licence metadata for this exact version and shown on the abstract page https://arxiv.org/abs/2508.08145v2. Full licence notice: \texttt{licenses/arxiv-2508.08145-CC-BY-4.0.txt}.\par\smallskip

\noindent\textit{Editorial scope.} Counted unit: the article's prop cuspscusp (source0.tex lines 240--252). The article's complete original proof of it is delivered with it (source0.tex lines 254--269, one \texttt{proof} environment of 16 lines). No mathematics is written, repaired or re-derived: the statement and the proof are the article's own lines, in the article's own notation, and no retained line is substituted or re-flowed.  Retained uncounted as the source-local context the proof needs: lines 110--117; lines 142--156.  Nothing was omitted from any retained span, so the package carries no omission record.  No retained line contains a citation, so the wrapper carries no bibliography and no bibliography entry.  No retained line contains a figure environment, an image-inclusion command or drawing code, so no image is delivered and the package declares no asset dependency.  Editorial formatting only, in addition: an AMS \texttt{amsart} preamble that restates the article's own declarations and macros used by the retained lines, namely the environments \texttt{lemma}, \texttt{prop} on separate counters, together with the standard packages those lines need.  The retained environments print in this document as Lemma 1, Proposition 1 in order of appearance, whereas the article prints the same items with the numbering of its own full document.  The complete native source of the article as distributed in the arXiv e-print for this exact version is bundled unchanged as \texttt{sources/183074/source0.tex}.
\begin{lemma}\label{chern}
We have
\begin{eqnarray*}
\overline c_1   &=&  \psi^*c_1'-c_1\\
\overline c_2 &=&  \psi^*c_2'-c_2-\psi^*c_1'c_1+c_1^2\\
\overline c_3  &=&  \psi^*c_3'-c_3-\psi^*c_2'c_1-\psi^*c_1'c_2+\psi^*c_1'c_1^2+2c_1c_2-c_1^3.
\end{eqnarray*}
\end{lemma}

\begin{lemma}\label{chern2}
We have
\begin{eqnarray*}
\psi^*c_i'&=&\binom{n+1}ic_1(\mathcal O_{\mathbb P(\mathcal F)}(1))^i\\
c_1&=&-\pi^*c_1(\Omega^1_X)-\pi^*c_1(\mathcal F)+(n-r+1)c_1(\mathcal O_{\mathbb P(\mathcal F)}(1))\\
c_2&=& \pi^*(c_2(\Omega^1_X)+c_2(\mathcal F)+c_1(\Omega^1_X)c_1(\mathcal F))
-(n-r+1)\pi^*c_1(\Omega^1_X)c_1(\mathcal O_{\mathbb P(\mathcal F)}(1))\\
&&-(n-r)\pi^*c_1(\mathcal F)c_1(\mathcal O_{\mathbb P(\mathcal F)}(1))
+\binom{n-r+1}2c_1(\mathcal O_{\mathbb P(\mathcal F)}(1))^2\\
c_3&=&-\sum_{i=0}^3(-1)^{3-i}\binom{n-r+1-i}{3-i}\pi^*c_i(\mathcal F)c_1(\mathcal O_{\mathbb P(\mathcal F)}(1))^{3-i} \\
&&-\pi^*c_1(\Omega_X^1)\sum_{i=0}^2(-1)^{2-i}\binom{n-r+1-i}{2-i}\pi^*c_i(\mathcal F)c_1(\mathcal O_{\mathbb P(\mathcal F)}(1))^{2-i} \\
&&-\pi^*c_2(\Omega_X^1)\sum_{i=0}^1(-1)^{1-i}\binom{n-r+1-i}{1-i}\pi^*c_i(\mathcal F)c_1(\mathcal O_{\mathbb P(\mathcal F)}(1))^{1-i} \\
&& -\pi^*c_3(\Omega_X^1).
\end{eqnarray*}
\end{lemma}

\begin{prop}\label{cuspscusp}
Set $h:=c_1(\mathcal O_{\mathbb P(\mathcal F)}(1))=\psi^*c_1(\mathcal O_{\mathbb P(V)}(1)$, $w_i:=\pi^*c_i(\Omega^1_X)$, and $f_i:=\pi^*c_i(\mathcal F)$. 
The class of the cuspidal locus $\kappa_\mathcal F:=\psi_*(\Sigma^{1,1,1})$ of the cuspidal edge $C_\mathcal F$ is
\begin{multline*}\textstyle
[\kappa_\mathcal F]=\psi_*\bigl(
(17\binom{r}3 +12\binom{r}2 +r)h^3
+\bigl((17\binom{r}2 + 6r)w_1
 \textstyle +  (17\binom{r}2 +r +2)f_1\bigr)h^2\\
 +\bigl(11rw_1^2+(17r-5)w_1f_1+(11r-7)f_1^2-5rw_2-(5r-4)f_2\bigr)h\\
 +6w_1^3+11w_1^2f_1+11w_1f_1^2+6f_1^3-7w_1w_2
 -7w_2f_1-5w_1f_2-7f_1f_2+2w_3+2f_3\bigr)\cap [\mathbb P(V)].
\end{multline*}
\end{prop}

\begin{proof} The class of $\kappa_\mathcal F$ is given by the Thom polynomial of $\Sigma^{1,1,1}$, which is equal to $\overline c_1^3+3\overline c_1\overline c_2+2\overline c_3$.  Lemma \ref{chern} and Lemma \ref{chern2} give 
\begin{multline*}
\textstyle \overline c_1=rc_1(\mathcal O_{\mathbb P(\mathcal F)}(1))+\pi^*(c_1(\Omega^1_X)+c_1(\mathcal F) )\\
\textstyle \overline c_2=\binom{r}2 c_1(\mathcal O_{\mathbb P(\mathcal F)}(1))^2+\pi^*(rc_1(\Omega^1_X)+(r-1)c_1(\mathcal F))c_1(\mathcal O_{\mathbb P(\mathcal F)}(1))\\
+\pi^*(c_1(\Omega^1_X)^2+c_1(\Omega^1_X)c_1(\mathcal F)+c_1(\mathcal F)^2
-c_2(\Omega^1_X)-c_2(\mathcal F))\\
\textstyle \overline c_3=\binom{r}3c_1(\mathcal O_{\mathbb P(\mathcal F)}(1))^3 +\pi^*(\binom{r}2 c_1(\Omega_X^1)+\binom{r-1}2 c_1(\mathcal F))c_1(\mathcal O_{\mathbb P(\mathcal F)}(1))^2 \\+\pi^*\bigl(rc_1(\Omega_X^1)^2+(r-1)c_1(\Omega_X^1)c_1(\mathcal F)
+(r-2)c_1(\mathcal F)^2-rc_2(\Omega_X^1)\\
-(r-2)c_2(\mathcal F)\bigr)c_1(\mathcal O_{\mathbb P(\mathcal F)}(1))
+\pi^*\bigl(c_3(\mathcal F)+c_3(\Omega_X^1)+c_1(\Omega_X^1)c_2(\mathcal F)\\
-2c_1(\Omega_X^1)c_2(\Omega_X^1)-2c_1(\Omega_X^1)c_2(\mathcal F)\-2c_2(\Omega_X^1)c_1(\mathcal F)
-2c_1(\mathcal F)c_2(\mathcal F)\\
+c_1(\Omega_X^1)^3+c_1(\Omega_X^1)^2c_1(\mathcal F)+c_1(\Omega_X^1)c_1(\mathcal F)^2+c_1(\mathcal F)^3\bigr),
\end{multline*}
from which the proposition follows. Note that the expression is polynomial in $r$ and independent of $n$.
\end{proof}

\end{document}
